% ══════════════════════════════════════════════════════════════════
% Section 2 — Continuum Mechanics Background
% ══════════════════════════════════════════════════════════════════
\section{Continuum mechanics framework}
\label{sec:continuum}

\subsection{Kinematics}
\label{sec:continuum:kinematics}

Consider a body undergoing a deformation described by the deformation gradient
$\Fmat = \pd{\tens{x}}{\tens{X}} \in \R^{3 \times 3}$, where $\tens{X}$ and
$\tens{x}$ denote material and spatial coordinates, respectively.  The right
Cauchy--Green deformation tensor is $\Cmat = \Fmat^T \Fmat$, with determinant
$J = \det \Fmat > 0$.

For nearly incompressible materials, the multiplicative decomposition
$\Fmat = J^{1/3}\,\bar{\Fmat}$ yields the isochoric part
$\bar{\Cmat} = J^{-2/3}\,\Cmat$ and the isochoric invariants:
\begin{align}
  \bar{I}_1 &= \tr(\bar{\Cmat}), &
  \bar{I}_2 &= \tfrac{1}{2}\bigl[\bar{I}_1^2 - \tr(\bar{\Cmat}^2)\bigr].
  \label{eq:isochoric_invariants}
\end{align}
For fiber-reinforced materials with preferred direction $\tens{a}_0$, the
pseudo-invariants are:
\begin{align}
  I_4 &= \tens{a}_0 \cdot \Cmat\,\tens{a}_0, &
  I_5 &= \tens{a}_0 \cdot \Cmat^2\,\tens{a}_0.
  \label{eq:fiber_invariants}
\end{align}

\subsection{Hyperelastic constitutive relations}
\label{sec:continuum:constitutive}

A hyperelastic material is characterized by a strain energy function (SEF)
$\sef$ such that the second Piola--Kirchhoff stress and the material tangent
are:
\begin{align}
  \Smat &= 2\,\pd{\sef}{\Cmat}, &
  \mathbb{C} &= 4\,\frac{\partial^2 \sef}{\partial \Cmat \, \partial \Cmat}.
  \label{eq:pk2_tangent}
\end{align}
Expressing $\sef$ in terms of invariants $\{I_\alpha\}$, the chain rule gives:
\begin{equation}
  \Smat = 2 \sum_\alpha \pd{\sef}{I_\alpha}\,\pd{I_\alpha}{\Cmat},
  \label{eq:pk2_chain}
\end{equation}
where the invariant derivatives $\partial I_\alpha / \partial \Cmat$ are known
analytically.  This decomposition is central to the hybrid NN--UMAT approach
(\Cref{sec:fortran_export}): the neural network learns
$\sef(I_1, I_2, \ldots)$ and its invariant gradients $\partial \sef / \partial
I_\alpha$, while the kinematics remain analytical.

\subsection{Material models}
\label{sec:continuum:materials}

We consider five materials spanning isotropic and anisotropic behavior.

\subsubsection{Isotropic models}

\paragraph{Neo-Hookean.}
\begin{equation}
  \sef = C_{10}(\bar{I}_1 - 3) + \frac{K}{2}(J - 1)^2.
  \label{eq:neohooke}
\end{equation}

\paragraph{Mooney--Rivlin.}
\begin{equation}
  \sef = C_{10}(\bar{I}_1 - 3) + C_{01}(\bar{I}_2 - 3) + \frac{K}{2}(J - 1)^2.
  \label{eq:mooneyrivlin}
\end{equation}

\paragraph{Yeoh.}
\begin{equation}
  \sef = \sum_{i=1}^{3} C_{i0}(\bar{I}_1 - 3)^i + \frac{K}{2}(J - 1)^2.
  \label{eq:yeoh}
\end{equation}

\subsubsection{Anisotropic models}

\paragraph{Holzapfel--Ogden (HO).}
A fiber-reinforced model with one family of fibers ($\tens{a}_0$):
\begin{equation}
  \sef = \frac{a}{2b}\bigl[e^{b(\bar{I}_1 - 3)} - 1\bigr]
       + \frac{a_f}{2b_f}\bigl[e^{b_f\langle I_4 - 1 \rangle^2} - 1\bigr]
       + \frac{K}{2}(J - 1)^2,
  \label{eq:holzapfelogden}
\end{equation}
where $\langle \cdot \rangle$ denotes the Macaulay bracket (tension-only fiber
contribution).

\paragraph{Gasser--Ogden--Holzapfel (GOH).}
Extends HO with a fiber dispersion parameter $\kappa \in [0, 1/3]$:
\begin{equation}
  \sef = \frac{a}{2b}\bigl[e^{b(\bar{I}_1 - 3)} - 1\bigr]
       + \frac{a_f}{2b_f}\Bigl[e^{b_f\bigl[\kappa \bar{I}_1
         + (1-3\kappa)I_4 - 1\bigr]^2} - 1\Bigr]
       + \frac{K}{2}(J - 1)^2.
  \label{eq:goh}
\end{equation}

% ── Summary table ─────────────────────────────────────────────────
\begin{table}[htbp]
\centering
\caption{Material models, invariant inputs, and default parameters used in the
  benchmarks.}
\label{tab:materials}
\begin{tabular}{@{}llcl@{}}
\toprule
\textbf{Model} & \textbf{Type} & \textbf{Input dim} & \textbf{Parameters} \\
\midrule
Neo-Hookean    & Isotropic   & 3 ($\bar{I}_1, \bar{I}_2, J$)
  & $C_{10}=0.5$, $K=100$ \\
Mooney--Rivlin & Isotropic   & 3
  & $C_{10}=0.5$, $C_{01}=0.2$, $K=100$ \\
Yeoh           & Isotropic   & 3
  & $C_{10}=0.5$, $C_{20}=0.1$, $C_{30}=0.01$, $K=100$ \\
\midrule
Holzapfel--Ogden & Anisotropic & 5 ($\bar{I}_1, \bar{I}_2, J, I_4, I_5$)
  & $a=0.059$, $b=4$, $a_f=2$, $b_f=4$, $K=100$ \\
GOH              & Anisotropic & 5
  & as HO $+\;\kappa=0.226$ \\
\bottomrule
\end{tabular}
\end{table}
